\question[4]
Erkläre in ein bis zwei Sätzen.
\begin{parts}
\part Was ist ein Abstrakter Datentyp?
\part Warum zeigt \texttt{pos} auf den Eintrag vor dem aktuellen Element und nicht auf das aktuelle Element selbst?
\part Wozu dient der Dummy-Eintrag am Anfang der Liste?
\part Warum nennt man die Liste eine rekursive Datenstruktur?
\end{parts}
\begin{solutionbox}{5cm}
a) Ein ADT besteht aus einer Datenstruktur und Operationen darauf. Festgelegt ist nur die Schnittstelle, nicht die Implementierung. \\
b) Beim Einfügen und Löschen muss der next-Zeiger des Vorgängers geändert werden. Vom aktuellen Element aus kommen wir
nicht zurück zum Vorgänger. \\
c) Dadurch hat auch das erste Element einen Vorgänger. Einfügen und Löschen am Anfang brauchen keinen Sonderfall. \\
d) Die Klasse Eintrag hat mit next ein Attribut vom Typ Eintrag. Sie verweist also auf sich selbst (rekursive Assoziation).
\end{solutionbox}
